% Options for packages loaded elsewhere
\PassOptionsToPackage{unicode}{hyperref}
\PassOptionsToPackage{hyphens}{url}
\PassOptionsToPackage{dvipsnames,svgnames,x11names}{xcolor}
\documentclass[
  a4paper,
  11pt]{article}
\usepackage{xcolor}
\usepackage[margin=2.5cm]{geometry}
\usepackage{amsmath,amssymb}
\setcounter{secnumdepth}{-\maxdimen} % remove section numbering
\usepackage{iftex}
\ifPDFTeX
  \usepackage[T1]{fontenc}
  \usepackage[utf8]{inputenc}
  \usepackage{textcomp} % provide euro and other symbols
\else % if luatex or xetex
  \usepackage{unicode-math} % this also loads fontspec
  \defaultfontfeatures{Scale=MatchLowercase}
  \defaultfontfeatures[\rmfamily]{Ligatures=TeX,Scale=1}
\fi
\usepackage{lmodern}
\ifPDFTeX\else
  % xetex/luatex font selection
\fi
% Use upquote if available, for straight quotes in verbatim environments
\IfFileExists{upquote.sty}{\usepackage{upquote}}{}
\IfFileExists{microtype.sty}{% use microtype if available
  \usepackage[]{microtype}
  \UseMicrotypeSet[protrusion]{basicmath} % disable protrusion for tt fonts
}{}
\makeatletter
\@ifundefined{KOMAClassName}{% if non-KOMA class
  \IfFileExists{parskip.sty}{%
    \usepackage{parskip}
  }{% else
    \setlength{\parindent}{0pt}
    \setlength{\parskip}{6pt plus 2pt minus 1pt}}
}{% if KOMA class
  \KOMAoptions{parskip=half}}
\makeatother
\usepackage{color}
\usepackage{fancyvrb}
\newcommand{\VerbBar}{|}
\newcommand{\VERB}{\Verb[commandchars=\\\{\}]}
\DefineVerbatimEnvironment{Highlighting}{Verbatim}{commandchars=\\\{\}}
% Add ',fontsize=\small' for more characters per line
\newenvironment{Shaded}{}{}
\newcommand{\AlertTok}[1]{\textcolor[rgb]{1.00,0.00,0.00}{\textbf{#1}}}
\newcommand{\AnnotationTok}[1]{\textcolor[rgb]{0.38,0.63,0.69}{\textbf{\textit{#1}}}}
\newcommand{\AttributeTok}[1]{\textcolor[rgb]{0.49,0.56,0.16}{#1}}
\newcommand{\BaseNTok}[1]{\textcolor[rgb]{0.25,0.63,0.44}{#1}}
\newcommand{\BuiltInTok}[1]{\textcolor[rgb]{0.00,0.50,0.00}{#1}}
\newcommand{\CharTok}[1]{\textcolor[rgb]{0.25,0.44,0.63}{#1}}
\newcommand{\CommentTok}[1]{\textcolor[rgb]{0.38,0.63,0.69}{\textit{#1}}}
\newcommand{\CommentVarTok}[1]{\textcolor[rgb]{0.38,0.63,0.69}{\textbf{\textit{#1}}}}
\newcommand{\ConstantTok}[1]{\textcolor[rgb]{0.53,0.00,0.00}{#1}}
\newcommand{\ControlFlowTok}[1]{\textcolor[rgb]{0.00,0.44,0.13}{\textbf{#1}}}
\newcommand{\DataTypeTok}[1]{\textcolor[rgb]{0.56,0.13,0.00}{#1}}
\newcommand{\DecValTok}[1]{\textcolor[rgb]{0.25,0.63,0.44}{#1}}
\newcommand{\DocumentationTok}[1]{\textcolor[rgb]{0.73,0.13,0.13}{\textit{#1}}}
\newcommand{\ErrorTok}[1]{\textcolor[rgb]{1.00,0.00,0.00}{\textbf{#1}}}
\newcommand{\ExtensionTok}[1]{#1}
\newcommand{\FloatTok}[1]{\textcolor[rgb]{0.25,0.63,0.44}{#1}}
\newcommand{\FunctionTok}[1]{\textcolor[rgb]{0.02,0.16,0.49}{#1}}
\newcommand{\ImportTok}[1]{\textcolor[rgb]{0.00,0.50,0.00}{\textbf{#1}}}
\newcommand{\InformationTok}[1]{\textcolor[rgb]{0.38,0.63,0.69}{\textbf{\textit{#1}}}}
\newcommand{\KeywordTok}[1]{\textcolor[rgb]{0.00,0.44,0.13}{\textbf{#1}}}
\newcommand{\NormalTok}[1]{#1}
\newcommand{\OperatorTok}[1]{\textcolor[rgb]{0.40,0.40,0.40}{#1}}
\newcommand{\OtherTok}[1]{\textcolor[rgb]{0.00,0.44,0.13}{#1}}
\newcommand{\PreprocessorTok}[1]{\textcolor[rgb]{0.74,0.48,0.00}{#1}}
\newcommand{\RegionMarkerTok}[1]{#1}
\newcommand{\SpecialCharTok}[1]{\textcolor[rgb]{0.25,0.44,0.63}{#1}}
\newcommand{\SpecialStringTok}[1]{\textcolor[rgb]{0.73,0.40,0.53}{#1}}
\newcommand{\StringTok}[1]{\textcolor[rgb]{0.25,0.44,0.63}{#1}}
\newcommand{\VariableTok}[1]{\textcolor[rgb]{0.10,0.09,0.49}{#1}}
\newcommand{\VerbatimStringTok}[1]{\textcolor[rgb]{0.25,0.44,0.63}{#1}}
\newcommand{\WarningTok}[1]{\textcolor[rgb]{0.38,0.63,0.69}{\textbf{\textit{#1}}}}
\usepackage{longtable,booktabs,array}
\usepackage{calc} % for calculating minipage widths
% Correct order of tables after \paragraph or \subparagraph
\usepackage{etoolbox}
\makeatletter
\patchcmd\longtable{\par}{\if@noskipsec\mbox{}\fi\par}{}{}
\makeatother
% Allow footnotes in longtable head/foot
\IfFileExists{footnotehyper.sty}{\usepackage{footnotehyper}}{\usepackage{footnote}}
\makesavenoteenv{longtable}
\setlength{\emergencystretch}{3em} % prevent overfull lines
\providecommand{\tightlist}{%
  \setlength{\itemsep}{0pt}\setlength{\parskip}{0pt}}
\usepackage[]{biblatex}
\addbibresource{references.bib}
\usepackage{etoolbox}
\AtBeginEnvironment{longtable}{\footnotesize}
\setlength{\tabcolsep}{4pt}
\usepackage{bookmark}
\IfFileExists{xurl.sty}{\usepackage{xurl}}{} % add URL line breaks if available
\urlstyle{same}
\hypersetup{
  pdftitle={The primes and the fence: an attack on Firoozbakht\textquotesingle s conjecture},
  pdfauthor={Noogram},
  colorlinks=true,
  linkcolor={black},
  filecolor={Maroon},
  citecolor={MidnightBlue},
  urlcolor={MidnightBlue},
  pdfcreator={LaTeX via pandoc}}

\title{The primes and the fence: an attack on
Firoozbakht\textquotesingle s conjecture}
\author{Noogram}
\date{2026-07-24}

\begin{document}
\maketitle

\nocite{*}

\subsection{0. Posture, stated before anything
else}\label{0-posture-stated-before-anything-else}

\subsubsection{0.1 The conjecture is
open}\label{01-the-conjecture-is-open}

\begin{quote}
\textbf{Firoozbakht\textquotesingle s conjecture is open.} This paper
does \textbf{not} prove it and does \textbf{not} refute it. No leg of
the polymer that produced this paper generated any evidence about its
truth in either direction, and this paper generates none either.

Everything below is a report on \emph{arguments, machine-checked
formalisations, computations and citations}. Nowhere is a statement made
about the conjecture\textquotesingle s truth value beyond the single
word \textbf{open}.
\end{quote}

The word \textbf{proved} is used in this paper under one rule and no
other:

\begin{quote}
\textbf{Proof discipline.} A statement is called \textbf{proved} only
where a Lean 4 kernel checked it and \texttt{\#print\ axioms} showed the
clean line \texttt{{[}propext,\ Classical.choice,\ Quot.sound{]}} --- no
\texttt{sorryAx}, no \texttt{Lean.ofReduceBool}. Every such statement is
listed in §4 with its verbatim Lean type. Statements established by
human argument are called \textbf{established} or \textbf{argued}, never
proved, and each carries its weakest premise\textquotesingle s tier
(§1.3).
\end{quote}

\subsubsection{0.2 Delivery posture: staged, and
why}\label{02-delivery-posture-staged-and-why}

This paper is \textbf{staged, not cleared}. Three gates stand between it
and release, and none of them has passed:

\begin{longtable}[]{@{}lll@{}}
\toprule\noalign{}
gate & status & consequence for this paper \\
\midrule\noalign{}
\endhead
\bottomrule\noalign{}
\endlastfoot
\texttt{evidence-gate} & \textbf{BLOCKED} (§8) & two live BLOCKER faults
in the corpus\textquotesingle s self-description \\
\texttt{citation-gate} & \textbf{NOT RUN} & no citation in this paper
has been independently re-fetched and re-read by an auditor \\
\texttt{editorial-verdict} & \textbf{NOT RUN} & the adversarial review
of this document is a separate downstream molecule \\
\end{longtable}

§9 lists exactly what must close before this document can be released,
and §10 lists what it does not claim.

\subsubsection{0.3 Attribution}\label{03-attribution}

External attribution for this work is \textbf{Noogram}.
Firoozbakht\textquotesingle s conjecture is attributed to Faride
Firoozbakht (1982); §11.3 records that no primary document by
Firoozbakht exists, and states what may and may not be written about
that attribution.

\begin{center}\rule{0.5\linewidth}{0.5pt}\end{center}

\subsection{1. The problem, and the shape of
it}\label{1-the-problem-and-the-shape-of-it}

\subsubsection{1.1 Statement}\label{11-statement}

Let \(p_n\) denote the \(n\)-th prime, \textbf{1-indexed}: \(p_1 = 2\),
\(p_2 = 3\), \(p_3 = 5\), \ldots. Write \(g_n := p_{n+1} - p_n\) for the
\(n\)-th prime gap and \(L_n := \log p_n\) (natural logarithm
throughout).

\begin{quote}
\textbf{Firoozbakht\textquotesingle s conjecture (F).} For every
\(n \ge 1\), \[p_{n+1}^{\,1/(n+1)} \;<\; p_n^{\,1/n},\] equivalently:
the sequence \(\big(p_n^{1/n}\big)_{n \ge 1}\) is \textbf{strictly
decreasing}.
\end{quote}

The conjecture was proposed by Faride Firoozbakht in 1982 and remains
open. Its earliest public record is Rivera\textquotesingle s \emph{Prime
Puzzles} Conjecture 30 (2002) {[}\texttt{rivera2002conj30}, V1{]}; its
earliest citable print record is Ribenboim, \emph{The Little Book of
Bigger Primes}, 2nd ed., p. 185 {[}\texttt{ribenboim2004little}, V2{]}.

\subsubsection{1.2 Five forms, all pointwise
equivalent}\label{12-five-forms-all-pointwise-equivalent}

The conjecture admits five natural phrasings. They are \textbf{pointwise
equivalent} --- index by index, with no asymptotics and no exceptional
set. This is an elementary but load-bearing fact: it is what licenses
proving in one form and stating in another.

\begin{longtable}[]{@{}lll@{}}
\toprule\noalign{}
form & statement & ambient \\
\midrule\noalign{}
\endhead
\bottomrule\noalign{}
\endlastfoot
\textbf{F1} & \(p_{n+1}^{1/(n+1)} < p_n^{1/n}\) & \(\mathbb{R}\), real
power \\
\textbf{F2} & \(p_{n+1}^{\,n} < p_n^{\,n+1}\) & \(\mathbb{N}\), pure
integer \\
\textbf{F3} & \(n\log p_{n+1} < (n+1)\log p_n\) & \(\mathbb{R}\) \\
\textbf{F4} & \(a_{n+1} < a_n\), where \(a_n := (\log p_n)/n\) &
\(\mathbb{R}\) \\
\textbf{F5} & \(S_n > 0\), where
\(S_n := (n+1)\log p_n - n\log p_{n+1}\) & \(\mathbb{R}\) \\
\end{longtable}

\emph{Proof of equivalence.} \(p_n \ge 2 > 0\), so both sides of F1 are
positive and \(\log\) is strictly increasing: F1
\(\iff \tfrac{1}{n+1}\log p_{n+1} < \tfrac{1}{n}\log p_n\), and
multiplying by \(n(n+1) > 0\) gives F3. F3 \(\iff\) F2 because \(\exp\)
is strictly increasing and \(m^k = \exp(k \log m)\) for \(m \ge 1\). F3
\(\iff\) F4 by dividing by \(n(n+1) > 0\). F3 \(\iff\) F5 is the
definition of \(S_n\). \(\square\)

The bridges F1 \(\iff\) F2, F1 \(\iff\) F3 and F1 \(\iff\) (the gap form
below) are \textbf{proved} in the machine-checked sense of §0.1 --- see
§4.2.

\subsubsection{\texorpdfstring{1.3 The exact criterion: the conjecture
\emph{is} a statement about prime
gaps}{1.3 The exact criterion: the conjecture is a statement about prime gaps}}\label{13-the-exact-criterion-the-conjecture-is-a-statement-about-prime-gaps}

\begin{quote}
\textbf{Lemma 1 (exact criterion).} For every \(n \ge 1\),
\[\mathrm{F}\text{ holds at } n \quad\Longleftrightarrow\quad g_n \;<\; f_n,
\qquad f_n \;:=\; p_n^{\,1+1/n} - p_n \;=\; p_n\big(e^{L_n/n} - 1\big).\]
\end{quote}

\emph{Proof.} By §1.2, F at \(n\) \(\iff\)
\(n \log p_{n+1} < (n+1)\log p_n \iff \log p_{n+1} <
(1+\tfrac1n)L_n \iff p_{n+1} < p_n^{1+1/n}\); subtract \(p_n\). Every
step is an equivalence. \(\square\)

\textbf{There is no side condition, no asymptotic regime, no exceptional
small-\(n\) set, and no external explicit bound.} Lemma 1 is the whole
content of the conjecture, index by index. It is also \textbf{proved} in
the machine-checked sense (\texttt{firoozbakht\_iff\_gap}, §4.2), in the
equivalent formulation \(g_n < B_n\) with
\(B_n := p_n(p_n^{1/n}-1) = f_n\).

Call \(\sigma_n := f_n - g_n\) the \textbf{shortfall} at \(n\); F at
\(n\) \(\iff\) \(\sigma_n > 0\).

Two consequences organise everything that follows.

\begin{itemize}
\tightlist
\item
  \textbf{\(f_n\) is strictly decreasing in \(n\) for fixed \(p_n\).}
  Hence \emph{verifying} F needs an \textbf{upper} bound on the index
  \(n = \pi(p_n)\), and \emph{refuting} it needs a \textbf{lower} bound.
  Getting this backwards converts a verification into a refutation
  claim; it is the most common way to lose the thread in this subject.
\item
  \textbf{Asymptotically \(f_n = \log^2 p_n - \log p_n - 1 + o(1)\)}
  {[}\texttt{kourbatov2015upper} Thm 5, V0{]}. So the fence sits at
  height about \((\log p)^2\).
\end{itemize}

\subsubsection{1.4 The picture}\label{14-the-picture}

The primes walk along a road. Every step is a gap. Firoozbakht drew a
fence at height roughly \((\log p)^2\) and said: \textbf{the primes
never jump it --- not once, not ever, all the way to infinity.}

That picture makes both branches of the attack immediately legible.

\begin{itemize}
\tightlist
\item
  \textbf{To prove F}, one needs an upper bound on \emph{every} prime
  gap, of size \((\log p)^2\)-ish, with an explicit constant strictly
  below \(1\). Nothing in analytic number theory produces a
  polylogarithmic gap bound at all. The best unconditional bound is
  \(g_n \ll p_n^{0.525}\) {[}\texttt{bhp2001difference}, V2{]}; the
  Riemann hypothesis gives \(g_n \ll \sqrt{p_n}\log p_n\). Those are
  \textbf{powers of \(p\) against powers of \(\log p\)}. The distance is
  not quantitative --- it is categorical.
\item
  \textbf{To refute F}, one needs a single gap that clears the fence. At
  the record locus the gap reaches \(92.06\%\) of \((\log p)^2\) while
  the fence sits at \(97.06\%\) of the same quantity --- a shortfall
  factor of \(1.05426\) (§5.1). A few percent, in a quantity that took
  decades of distributed computation to measure.
\end{itemize}

The proof branch is blocked by a wall. The refutation branch is
\emph{close} but not reachable. §7.4 records how this framing was itself
demoted during the work.

\subsubsection{1.5 Tiering convention}\label{15-tiering-convention}

Two orthogonal ladders are used, and they are never conflated.

\textbf{Claim tier} --- how well a \emph{claim} is established:
\textbf{L0} proved here from first principles; \textbf{L1}
standard/textbook; \textbf{L2} attributed with a verified locator;
\textbf{L3} heuristic, recalled, or resting on an unrefereed
computational record; \textbf{{[}C{]}} additionally recomputed by a
verifier script in the run. A composite claim inherits its
\textbf{weakest} premise\textquotesingle s tier.

\textbf{Citation tier} --- how well a \emph{citation} was checked:
\textbf{V0} primary obtained and the statement read verbatim at the
named locator; \textbf{V1} primary\textquotesingle s own metadata page
read; \textbf{V2} confirmed via a reliable secondary quoting the
primary; \textbf{V3} bibliographic metadata only, statement at locator
\textbf{not} verified. Every citation in this paper is tagged, and §11
traces each to its source-ledger row.

\textbf{A V3 locator is not usable.} Where this paper touches one, it
says so and does not lean on it.

\begin{center}\rule{0.5\linewidth}{0.5pt}\end{center}

\subsection{2. Method}\label{2-method}

\subsubsection{2.1 What was actually
done}\label{21-what-was-actually-done}

The attack ran as a directed polymer of specialised legs, each producing
an artifact another leg had to read. In rough order:

\begin{enumerate}
\def\labelenumi{\arabic{enumi}.}
\tightlist
\item
  \textbf{Decomposition} --- reduce the conjecture to its equivalent
  forms and to the gap criterion (§1.2, §1.3); enumerate attack routes.
\item
  \textbf{Frame deliberation} --- a five-perspective panel adjudicating
  which branch is live.
\item
  \textbf{Source ledger} --- build the bibliography as \emph{(citekey,
  verification level, precise locator, exact statement as the source
  writes it)}, with independent recomputation of every numeric row.
\item
  \textbf{Concept cards} --- 31 atomic definitions/lemmas, each carrying
  its ledger row.
\item
  \textbf{Lean skeleton, then Lean probe} --- state the conjecture in
  Lean 4 / Mathlib and attempt real proof terms.
\item
  \textbf{Proof attempts} --- three targets, attacked as mathematics:
  first-failure maximality, the RH-conditional route, and the
  unconditional verified range.
\item
  \textbf{Notebooks} --- three computational legs, exhaustive and
  lever-based.
\item
  \textbf{Red-team corpus} --- 19 adversarial artifacts designed to slip
  past a naive gate.
\item
  \textbf{Skeptic} --- an adversarial audit that re-derived every
  load-bearing number in independent code and re-fetched five primary
  sources.
\item
  \textbf{Evidence gate} --- fail-closed adjudication.
\item
  \textbf{Synthesis}, then this paper.
\end{enumerate}

\subsubsection{2.2 Method commitments}\label{22-method-commitments}

Four commitments constrain everything reported here.

\begin{itemize}
\tightlist
\item
  \textbf{Fail-closed gates.} A gate that cannot establish its predicate
  returns FAIL, never "probably fine". The evidence gate returned
  \textbf{BLOCKED}, and §8 reports it as such.
\item
  \textbf{The kernel is the only source of the word "proved".} §0.1.
\item
  \textbf{The \texttt{sorry} grep is not the audit.} A raw
  \texttt{grep\ -c\ sorry} over the statement file returns 6, of which
  exactly \textbf{one} is a genuine \texttt{sorry} in code --- the rest
  are prose discussing what remains open. Verdicts rest on the
  \texttt{\#print\ axioms} whitelist, which cannot be fooled by
  formatting.
\item
  \textbf{Computation corroborates or refutes; it never proves.} No
  finite range bears on the conjecture. A certified range is a
  \emph{fidelity instrument}, not evidence.
\end{itemize}

\subsubsection{2.3 The formal
environment}\label{23-the-formal-environment}

Lean 4 \texttt{v4.29.0} with Mathlib \texttt{v4.29.0} (revision pinned
in the run\textquotesingle s \texttt{lake-manifest.json}).
\texttt{lake\ build} completes with \textbf{exit status 0}, 8253 jobs.
An \texttt{Audit.lean} module is part of the library, so
\texttt{\#print\ axioms} runs on \emph{every} build and cannot silently
rot. No \texttt{native\_decide} appears anywhere in the sources, so
\texttt{Lean.ofReduceBool} never enters the trusted base: the finite
verification is kernel-checked, not compiler-trusted.

\subsubsection{2.4 A named self-check on every result
claim}\label{24-a-named-self-check-on-every-result-claim}

The run\textquotesingle s red-team corpus produced an artifact
(\texttt{RT-12}) that builds with exit 0, is \texttt{sorry}-free, is
axiom-clean, \textbf{and proves a true theorem} --- where the falsehood
lives in the sentence a paper would put next to it. That artifact
imposes a standing requirement on this document:

\begin{quote}
\textbf{RT-12 requirement.} For every declaration this paper names as a
result, print its Lean type and confirm the conclusion is not among its
hypotheses.
\end{quote}

§4.4 discharges that requirement for all fourteen declarations named in
§4.

\begin{center}\rule{0.5\linewidth}{0.5pt}\end{center}

\subsection{3. Results at a glance}\label{3-results-at-a-glance}

\begin{longtable}[]{@{}ll@{}}
\toprule\noalign{}
question & answer \\
\midrule\noalign{}
\endhead
\bottomrule\noalign{}
\endlastfoot
Is the conjecture \textbf{proved}? & \textbf{No.} The Lean declaration
\texttt{Firoozbakht.firoozbakht} still ends in \texttt{sorry}; its axiom
line contains \texttt{sorryAx}. \\
Is the conjecture \textbf{refuted}? & \textbf{No.} No counterexample at
any scale examined. \\
Was anything \textbf{proved} (kernel sense)? & \textbf{Yes} --- 11 new
theorems plus a 26-entry prime table, including the conjecture itself
for all \(1 \le n \le 25\) and two exact refutation criteria. §4. \\
Was anything \textbf{established on paper}? & \textbf{Yes} --- §5. Two
exact criteria, a verified-range theorem at \(H = 10^{20}\) under two
named premises, and three dead routes with quantitative obituaries.
§6. \\
Evidence-gate status & \textbf{BLOCKED}. §8. \\
Citation clearance & \textbf{Not claimed, not audited.} §9. \\
\end{longtable}

\begin{center}\rule{0.5\linewidth}{0.5pt}\end{center}

\subsection{4. What was PROVED --- machine-checked in Lean 4 /
Mathlib}\label{4-what-was-proved--machine-checked-in-lean-4--mathlib}

Everything in this section satisfies §0.1: \texttt{lake\ build} exit 0,
and \texttt{\#print\ axioms} showing
\texttt{{[}propext,\ Classical.choice,\ Quot.sound{]}} for each
declaration.

Throughout, \texttt{nthPrime\ n} is the Lean 1-indexed \(n\)-th prime,
\texttt{noncomputable\ def\ nthPrime\ (n\ :\ ℕ)\ :\ ℕ\ :=\ Nat.nth\ Nat.Prime\ (n\ -\ 1)},
and \texttt{barrier\ n} is the exact barrier \(B_n = p_n(p_n^{1/n}-1)\)
of Lemma 1.

\subsubsection{4.1 The headline: the conjecture, proved outright, for 25
consecutive
indices}\label{41-the-headline-the-conjecture-proved-outright-for-25-consecutive-indices}

\begin{Shaded}
\begin{Highlighting}[]
\NormalTok{theorem firoozbakht\_of\_le (n : ℕ) (h1 : 1 ≤ n) (h2 : n ≤ 25) :}
\NormalTok{    (nthPrime (n + 1) : ℝ) \^{} ((1 : ℝ) / ((n : ℝ) + 1)) \textless{}}
\NormalTok{      (nthPrime n : ℝ) \^{} ((1 : ℝ) / (n : ℝ))}
\end{Highlighting}
\end{Shaded}

This is the \textbf{literal F1 form} --- real powers, not the integer
surrogate --- and its axiom line is clean:

\begin{verbatim}
'Firoozbakht.firoozbakht_of_le' depends on axioms: [propext, Classical.choice, Quot.sound]
\end{verbatim}

so it is \textbf{not circular}: it does not depend on the
\texttt{sorry}-ed conjecture. It is a genuine unconditional fragment of
an open problem, checked by the same kernel that checks the bridges.

\textbf{Honest sizing, stated immediately and without softening.} The
bound \(n \le 25\) is set by \textbf{compile time, not by mathematics}
--- the bridge lemma is uniform in \(n\) and the table extends
mechanically. And \(n \le 25\) is derisory next to the
literature\textquotesingle s \(4\times10^{18}\)
{[}\texttt{kourbatov2015verification}{]} and \(2^{64}\)
{[}\texttt{visser2019verifying}{]}. \textbf{No finite range bears on the
conjecture.} The value of this result is not its range; it is that (a)
the formal method now exists, and (b) a kernel-checked range is a
different kind of object from a script-checked one.

\subsubsection{\texorpdfstring{4.2 The enabling engineering: a
computable bridge out of
\texttt{Nat.nth}}{4.2 The enabling engineering: a computable bridge out of Nat.nth}}\label{42-the-enabling-engineering-a-computable-bridge-out-of-natnth}

\texttt{Nat.nth\ Nat.Prime} is \texttt{noncomputable} --- defined by
well-founded recursion over an infinite set --- so \texttt{decide}
cannot evaluate it and \textbf{no finite range could be machine-checked
at all}. This was the formalisation\textquotesingle s flagged main risk
item.

The escape hatch is
\texttt{Nat.nth\_count\ :\ p\ n\ →\ Nat.nth\ p\ (Nat.count\ p\ n)\ =\ n},
which trades \texttt{nth} for the computable \texttt{count}. The naive
move works but is quadratic in the wrong variable: a single check
\texttt{Nat.count\ Nat.Prime\ 97\ =\ 24} walks the whole range
\([0,97)\) and cost \textbf{\textasciitilde36 s} of kernel reduction as
measured (one entry, the largest prime in the table). That does not
scale.

The leg instead proves a step lemma whose cost scales with the
\textbf{gap}, not the prime:

\begin{Shaded}
\begin{Highlighting}[]
\NormalTok{theorem nthPrime\_succ\_eq \{k p q : ℕ\} (hk : 1 ≤ k) (hp : nthPrime k = p) (hq : Nat.Prime q)}
\NormalTok{    (hpq : p \textless{} q) (hgap : ∀ m, p \textless{} m → m \textless{} q → ¬ Nat.Prime m) :}
\NormalTok{    nthPrime (k + 1) = q}
\end{Highlighting}
\end{Shaded}

Each \texttt{hgap} obligation is an \texttt{interval\_cases} over a gap
of size \(< 10\) in this range --- milliseconds. The whole file (both
bridge lemmas, the 26-entry table, \emph{and} the certified range)
compiles in \textbf{6.5--9.6 s}. \textbf{No numeral is asserted
anywhere}: \(p_1 = 2\) comes from Mathlib, and every later entry is
\emph{derived} from its predecessor.

Also proved and used: \texttt{count\_prime\_eq\_of\_no\_prime} (no prime
in \([a,b)\) \(\Rightarrow\) \(\pi\) constant across it),
\texttt{nthPrime\_succ\_le\_of\_prime} (minimality), and the equivalence
bridges \texttt{firoozbakht\_iff\_log}, \texttt{firoozbakht\_iff\_int},
\texttt{firoozbakht\_iff\_gap} --- the machine-checked form of §1.2 and
Lemma 1.

\subsubsection{4.3 Both directional criteria,
exact}\label{43-both-directional-criteria-exact}

\textbf{Sufficient (verification side).}

\begin{Shaded}
\begin{Highlighting}[]
\NormalTok{theorem firoozbakht\_of\_gap\_lt \{n : ℕ\} (hn : 1 ≤ n)}
\NormalTok{    (h : (n : ℝ) * ((nthPrime (n + 1) : ℝ) {-} (nthPrime n : ℝ))}
\NormalTok{        \textless{} (nthPrime n : ℝ) * Real.log (nthPrime n)) :}
\NormalTok{    (nthPrime (n + 1) : ℝ) \^{} ((1 : ℝ) / ((n : ℝ) + 1)) \textless{}}
\NormalTok{      (nthPrime n : ℝ) \^{} ((1 : ℝ) / (n : ℝ))}
\end{Highlighting}
\end{Shaded}

i.e. \(n g_n < p_n \log p_n \Rightarrow \mathrm{F}\) at \(n\). The proof
is the single estimate \(\log(p_{n+1}/p_n) \le p_{n+1}/p_n - 1\), which
is exactly where the exact barrier degrades to its first-order
surrogate. The hypothesis is therefore \emph{strictly stronger} than the
conjecture at \(n\): \textbf{sufficient, never necessary.}

\textbf{Necessary conditions for a refutation --- two exact criteria.}

\begin{Shaded}
\begin{Highlighting}[]
\NormalTok{theorem not\_firoozbakht\_of\_barrier\_le \{n : ℕ\} (hn : 1 ≤ n)}
\NormalTok{    (h : barrier n ≤ ((nthPrime (n + 1) : ℝ) {-} (nthPrime n : ℝ))) :}
\NormalTok{    ¬ ((nthPrime (n + 1) : ℝ) \^{} ((1 : ℝ) / ((n : ℝ) + 1)) \textless{}}
\NormalTok{        (nthPrime n : ℝ) \^{} ((1 : ℝ) / (n : ℝ)))}

\NormalTok{theorem not\_firoozbakht\_of\_pow\_le \{n : ℕ\} (hn : 1 ≤ n)}
\NormalTok{    (h : (nthPrime n) \^{} (n + 1) ≤ (nthPrime (n + 1)) \^{} n) :}
\NormalTok{    ¬ ((nthPrime (n + 1) : ℝ) \^{} ((1 : ℝ) / ((n : ℝ) + 1)) \textless{}}
\NormalTok{        (nthPrime n : ℝ) \^{} ((1 : ℝ) / (n : ℝ)))}
\end{Highlighting}
\end{Shaded}

The second is \textbf{pure \(\mathbb{N}\) arithmetic} --- no real
arithmetic, no \texttt{Real.log} --- and is therefore the right target
for a computational search. Note what it demands: the \textbf{index}
\(n = \pi(p_n)\) must be certified, not merely the gap. \emph{A large
gap alone is not a counterexample.}

\textbf{Explicitly not admissible as a refutation}: exhibiting a gap
above \(\log^2 p_n - \log p_n - 1\). That is the asymptotic surrogate,
not \(B_n\), and the difference is \(O(1/\log p_n)\) --- which is
exactly the size of the margin being contested.

\subsubsection{4.4 RT-12 discharge --- hypothesis/conclusion
separation}\label{44-rt-12-discharge--hypothesisconclusion-separation}

Per §2.4, for every declaration this paper names as a result:

\begin{longtable}[]{@{}llll@{}}
\toprule\noalign{}
declaration & hypotheses & conclusion & conclusion among hypotheses? \\
\midrule\noalign{}
\endhead
\bottomrule\noalign{}
\endlastfoot
\texttt{firoozbakht\_of\_le} & \(1\le n\), \(n\le 25\) & F1 at \(n\) &
\textbf{no} \\
\texttt{firoozbakht\_gap\_of\_le} & \(1\le n\), \(n\le 25\) &
\(g_n < B_n\) & \textbf{no} \\
\texttt{nthPrime\_succ\_eq} & \(1\le k\), \texttt{nthPrime\ k\ =\ p},
\texttt{q} prime, \(p<q\), no prime strictly between &
\texttt{nthPrime\ (k+1)\ =\ q} & \textbf{no} \\
\texttt{nthPrime\_eq\_1\ …\ \_26} & none & \$p\_i = \$ numeral &
\textbf{no} \\
\texttt{count\_prime\_eq\_of\_no\_prime} & \(a\le b\), no prime in
\([a,b)\) & \(\pi\) constant & \textbf{no} \\
\texttt{nthPrime\_succ\_le\_of\_prime} & \(1\le n\), \texttt{q} prime,
\(p_n<q\) & \(p_{n+1}\le q\) & \textbf{no} \\
\texttt{firoozbakht\_iff\_log} & \(1\le n\) & F1 \(\iff\) F3 &
\textbf{no} (biconditional bridge) \\
\texttt{firoozbakht\_iff\_int} & \(1\le n\) & F1 \(\iff\) F2 &
\textbf{no} \\
\texttt{firoozbakht\_iff\_gap} & \(1\le n\) & F1 \(\iff\) \(g_n<B_n\) &
\textbf{no} \\
\texttt{firoozbakht\_of\_gap\_lt} & \(1\le n\), \(n g_n < p_n L_n\) & F1
at \(n\) & \textbf{no} \\
\texttt{not\_firoozbakht\_of\_barrier\_le} & \(1\le n\), \(B_n \le g_n\)
& \(\neg\)F1 at \(n\) & \textbf{no} \\
\texttt{not\_firoozbakht\_of\_pow\_le} & \(1\le n\),
\(p_n^{n+1}\le p_{n+1}^n\) & \(\neg\)F1 at \(n\) & \textbf{no} \\
\texttt{nthPrime\_succ\_lt\_two\_mul} & \(1\le n\) & \(p_{n+1}<2p_n\) &
\textbf{no} \\
\texttt{firoozbakht\_of\_two\_pow\_le} & \(1\le n\), \(2^n \le p_n\) &
F1 at \(n\) & \textbf{no} \\
\end{longtable}

No declaration in this table takes Firoozbakht\textquotesingle s
conjecture --- in any of its five forms --- as a hypothesis. The two
refutation criteria are genuine contrapositives; this was independently
confirmed inside the kernel by typechecking

\begin{Shaded}
\begin{Highlighting}[]
\NormalTok{example (n : ℕ) (h1 : 1 ≤ n) (h2 : n ≤ 25)}
\NormalTok{    (hbad : (nthPrime n) \^{} (n + 1) ≤ (nthPrime (n + 1)) \^{} n) : False :=}
\NormalTok{  not\_firoozbakht\_of\_pow\_le h1 hbad (firoozbakht\_of\_le n h1 h2)}
\end{Highlighting}
\end{Shaded}

which wires the refutation criterion and the certified range together
and shows the kernel agrees they are incompatible.

\subsubsection{4.5 An external cross-check on the
numerals}\label{45-an-external-cross-check-on-the-numerals}

Lean checks that the proofs follow. It does not check that the
\emph{numerals} are the primes anyone means. Against an independent
\texttt{sympy} sieve:

\begin{verbatim}
prime table p_1..p_26 matches sympy sieve: True
n in 1..25 failing p_{n+1}^n < p_n^(n+1): []
n with 2^n <= p_n (n < 40): [1]
n < 72849 where linearised criterion n*g_n < p_n*log p_n FAILS: count=2 first=[2, 4]
\end{verbatim}

The third line is a genuinely useful fact: over the first
\textasciitilde73 000 indices the linearised sufficient criterion fails
at only \(n = 2\) and \(n = 4\), \textbf{both inside the range already
proved outright} by \texttt{firoozbakht\_of\_le}. So
\texttt{firoozbakht\_of\_gap\_lt} is not a lossy detour --- and equally,
it is \textbf{no easier than the target}.

\begin{center}\rule{0.5\linewidth}{0.5pt}\end{center}

\subsection{5. What was ESTABLISHED on
paper}\label{5-what-was-established-on-paper}

These are arguments, not certificates. Each carries its tier per §1.5.

\subsubsection{5.1 The record locus, and the size of the
margin}\label{51-the-record-locus-and-the-size-of-the-margin}

At the largest known maximal-gap locus {[}\texttt{kourbatov2015upper}
Table 1, last row, \textbf{V0}{]}:

\[k = 49\,749\,629\,143\,526, \qquad p_k = 1\,693\,182\,318\,746\,371, \qquad g_k = 1132.\]

Recomputed here in 60-digit arithmetic from exact integer inputs
(\textbf{{[}C{]}}):

\begin{longtable}[]{@{}llll@{}}
\toprule\noalign{}
quantity & recomputed & published & source \\
\midrule\noalign{}
\endhead
\bottomrule\noalign{}
\endlastfoot
\(\log p_k\) & \(35.0653861820133348\ldots\) & --- & --- \\
\(f_k = p_k^{1+1/k}-p_k\) & \(1193.41777829404\ldots\) & \(1193.418\) &
\texttt{kourbatov2015upper} Table 1 ✓ \\
\(\ell_k = \log^2 p_k-\log p_k\) & \(1194.51592191172\ldots\) &
\(1194.516\) & \texttt{kourbatov2015upper} Table 1 ✓ \\
\(R_{\mathrm{csg}} = g_k/\log^2 p_k\) & \(0.9206385885574205\ldots\) &
\(0.9206\) & \texttt{primegaplist\_faq} ✓ \\
\(f_k/\log^2 p_k\) (the fence, same normalisation) &
\(0.9705887446713\ldots\) & --- & \textbf{{[}C{]}} \\
\textbf{shortfall \(f_k/g_k\)} & \(\mathbf{1.05425598789\ldots}\) & ---
& \textbf{{[}C{]}} \\
\end{longtable}

\textbf{Use \(1.05426\)}, sourced directly to the \emph{published exact
barrier} \(f_k = 1193.418\). Not \(1.055\) (an earlier figure in this
run, superseded), and not via any \(c_n\) interval --- the exact barrier
makes the explicit-bounds detour unnecessary. If an explicit-bounds
route is wanted anyway it certifies \(1.05424 \pm 0.00005\) (Dusart 2010
Props. 6.6/6.7, \textbf{V0}, validity \(k \ge 688\,383\)) ---
consistent, with the exact value inside it.

Tier: \textbf{L0/L1 + V0 + {[}C{]}}.

\textbf{A caveat that is real, not decorative.} The statement "the CSG
ratio \(R = g/\log^2 p\) has never exceeded 1" is \textbf{false as
written}. It requires the restriction \(p > 7\). Recomputed here:
\(R > 1\) at exactly three points --- \(p=2\) (\(R = 2.081369\)),
\(p=3\) (\(1.657071\)), \(p=7\) (\(1.056366\)) --- and \(R \le 1\) for
every other \(p \le 113\). The community source states the record
\emph{with} the restriction ("the largest value observed for any
\(p_1 > 7\)"); any restatement must carry it.

\subsubsection{\texorpdfstring{5.2 Theorem A --- an unconditional
verified range at
\(H = 10^{20}\)}{5.2 Theorem A --- an unconditional verified range at H = 10\^{}\{20\}}}\label{52-theorem-a--an-unconditional-verified-range-at-h--1020}

\begin{quote}
\textbf{Theorem A (established, not proved).} F holds for every \(n\)
with \(p_n < 10^{20}\).
\end{quote}

The argument is complete and every step is either first-principles or
one of exactly two external premises, both named and isolated:

\begin{longtable}[]{@{}lll@{}}
\toprule\noalign{}
premise & supplies & tier \\
\midrule\noalign{}
\endhead
\bottomrule\noalign{}
\endlastfoot
\textbf{P1} --- Axler, \emph{New bounds for the prime counting
function}, the \(3.83\) member:
\(\pi(x) < x\big/\big(\log x - 1 - \tfrac1{\log x} - \tfrac{3.83}{\log^2 x}\big)\)
for \(x \ge 9.25\) & the analytic criterion & \textbf{V0} · \textbf{L2}
--- see the locator fault below \\
\textbf{P2} --- the census of first-occurrence prime gaps is complete
below \(10^{20}\), yielding the 85 known maximal gaps & the range &
\textbf{V1} {[}\texttt{primegaplist\_faq}{]} · \textbf{L3}: a community
computational record, unrefereed, not machine-checked \\
\end{longtable}

Everything else --- the criterion, its monotonicity, the covering lemma,
the base case, the assembly --- is L0. By the propagation rule the
conclusion inherits its weakest premise: \textbf{Theorem A is L3/V1, and
its weak link is computational data, not mathematics.}

The structural device is a \textbf{lever}. If the maximal-gap table is
complete below \(H\), with records \((P_i, G_i)\), then for
\(P_i \le p_n < P_{i+1}\) we have \(g_n \le G_i\) --- a larger gap would
be an unrecorded new record. Combined with a lower bound on \(B_n\)
(from an \emph{upper} bound on \(\pi\), per §1.3), this turns
\textasciitilde80 inequalities into a verified range of
\(\approx 2.2\times10^{18}\) primes: \textbf{\(2.7\times10^{16}\) primes
per inequality} {[}COMPUTED{]}. That ratio \emph{is} the content of the
phrase "verified up to \(4\times10^{18}\)". Nobody has evaluated the
criterion at \(10^{17}\) consecutive primes, and nobody will.

Two subsidiary findings:

\begin{itemize}
\tightlist
\item
  \textbf{The range is data-bound, not method-bound.} At the tightest of
  the 85 records the criterion is satisfied with \(5.1\%\) to spare
  {[}C{]}; \(H\) stops exactly where the census stops. The gain from
  \(2^{64}\) (2019) to \(10^{20}\) (today) is \textbf{not mathematical}
  --- the same argument gives both; the census moved.
\item
  \textbf{The relaxation costs nothing.} Using the unconditional
  criterion instead of the exact barrier loses \(9.3\times10^{-6}\)
  relative {[}C{]}. This is why the target closes at all.
\end{itemize}

\textbf{A live citation fault on P1 --- MAJOR-6, stated here rather than
buried.} The Axler result is cited throughout this corpus as
\textbf{"Corollary 3.5"}, which is the \textbf{arXiv v3} numbering,
against a bibliographic entry whose journal-of-record fields are
\emph{Integers} \textbf{16} (2016), A22. The published corrigendum
numbers the affected result \textbf{Corollary 3.4, page 8}. A referee
checking the version of record will \textbf{not} find the \(3.83\) bound
at "Corollary 3.5". The premise is \emph{believed} --- the skeptic leg
retrieved arXiv v3 and read the four-member family verbatim, confirming
P1 at V0 --- but the locator does not resolve in the version it is
attributed to, and the corpus never establishes where the \(3.83\)
member sits in the published numbering. \textbf{Theorem A must not be
published until this locator is repaired.}

\subsubsection{5.3 The exhaustive floor the project owns
outright}\label{53-the-exhaustive-floor-the-project-owns-outright}

Independently of any external premise: F verified exhaustively for all
\(p < 10^9\) --- \textbf{50 847 533 consecutive pairs} --- by sieve plus
a \emph{runtime-certified} floating-point error bound, in 9 s. Tier
\textbf{L0 + {[}C{]}}. This is the only part of the range this work owns
without inheriting anyone\textquotesingle s computational claim. Below
\(p = 1223\) the check is exact integer comparison of \(p_{n+1}^n\)
against \(p_n^{n+1}\), with no arithmetic model at all.

Within that same floor, \textbf{FFM} --- \emph{if F first fails at
\(n^\star\), then \(g_{n^\star}\) is a record gap} --- is
\textbf{certified, not assumed}, for all \(n \le 50\,847\,533\).

\textbf{A vacuity note, applied here rather than inherited.} On any
range where F is \emph{verified}, FFM is \textbf{vacuously true} and
carries zero information: its antecedent is the negation of F. This
paper therefore does \textbf{not} advertise "FFM proved on
\([89, 10^{20})\)" as a deliverable, because on exactly that range F is
verified from exactly those premises. What is non-vacuous is the
\textbf{certificate} \(G_{n-1} < f_n\) on \([89, H)\) for every \(H\) to
which the census is complete --- strictly stronger than FFM-at-\(n\) and
not implied by F\textquotesingle s verification. The general FFM
statement, for unbounded \(n\), remains \textbf{open}.

\subsubsection{5.4 The FFM dichotomy}\label{54-the-ffm-dichotomy}

The real deliverable of the FFM target is a structural theorem, not a
range:

\begin{quote}
\textbf{Theorem B (established).} Every known route to proving FFM
passes either through F itself or through a statement strictly stronger
than F.
\end{quote}

Two elementary facts drive it. (a) F \(\Rightarrow\) FFM vacuously. (b)
Hence any \emph{refutation} of FFM contains a refutation of F --- to
exhibit a counterexample to FFM one must exhibit the least index at
which F fails, and in particular an index at which F fails. So FFM is
\emph{logically weaker} than F, \textbf{and the weakness buys nothing}.

Tier \textbf{L0}.

\begin{center}\rule{0.5\linewidth}{0.5pt}\end{center}

\subsection{6. What was REFUTED}\label{6-what-was-refuted}

Nothing about the conjecture. \textbf{Three routes} died, and the
obituaries are quantitative.

\subsubsection{6.1 The Riemann hypothesis route --- dead, with
numbers}\label{61-the-riemann-hypothesis-route--dead-with-numbers}

\begin{quote}
\textbf{Theorem C (established).} The strongest \emph{published}
RH-conditional prime-gap bound (Carneiro--Milinovich--Soundararajan
2019, constant \(22/25\)) decides Firoozbakht at \textbf{exactly three
indices}, \(n \in \{1,2,3\}\) --- and at \textbf{exactly one}, \(n = 3\)
(\(p = 5\)), once the source\textquotesingle s own validity restriction
\(p_n > 3\) is honoured. It is silent at every \(n \ge 4\), forever.
\end{quote}

\begin{quote}
\textbf{Theorem D (established).} This is not a constant problem. For
\textbf{every} \(c > 0\), a bound \(g_n \le c\sqrt{p_n}\log p_n\)
decides F on a finite initial segment only. Improving the constant moves
a threshold; it never changes the verdict.
\end{quote}

\begin{quote}
\textbf{Theorem E (established).} The RH branch is \textbf{strictly
dominated}. For every \(c \ge 10^{-7}\) --- seven orders of magnitude
below the published constant, and six below \(c = 1/2\), which the CMS
authors name as the limit of their method --- the usable set ends
\emph{inside} the range already verified \textbf{unconditionally}. RH
contributes \textbf{zero indices} that unconditional verification does
not already own.
\end{quote}

Tier \textbf{L0 + {[}C{]}} for all three, resting on a \textbf{V0}
reading of the CMS paper (but see the ledger caveat in §11.2).

The \(\sqrt{x}\) shape is not laziness. Littlewood\textquotesingle s
unconditional \(\Omega_\pm\) theorem forces any argument that bounds the
prime-counting error pointwise and subtracts to work at that scale
\textbf{regardless of RH}. (This explanatory proposition rests on a
\textbf{V3} row for Littlewood 1914 and is stated as such; Theorems C--E
do not use it.)

\textbf{What was deliberately \emph{not} proved}: that RH does not
\emph{imply} Firoozbakht. That is a claim about RH\textquotesingle s
deductive closure and nothing here establishes it. The barrier proved is
about the \textbf{shape of the bounds the known machinery emits}, over
an explicitly enumerated surveyed set. That restraint is deliberate and
is preserved in this paper.

\textbf{One incidental positive.} RH \emph{does} have a use in this
attack --- on the \textbf{refutation} branch, certifying the index
\(n = \pi(p)\). It is on the wrong branch of the tree, and the
unconditional bound was already comfortable there by a factor of
\textasciitilde1600.

\subsubsection{6.2 The monotone-barrier route to FFM --- dead, by an
exact
witness}\label{62-the-monotone-barrier-route-to-ffm--dead-by-an-exact-witness}

The tempting proof of FFM routes through monotonicity of the exact
barrier \(f_n\). \textbf{\(f_n\) is not monotone}: it decreases at
\(57.88\%\) of steps and sits below its own running maximum at
\(87.88\%\) of indices, with an \textbf{exact-integer witness at
\(n = 7, 8\)} (\(f_7 - f_8 = 0.0281326864429\), reproduced independently
in 40-digit arithmetic).

The error is subtle and worth recording precisely: the \textbf{proxy}
barrier \(\ell_n = \log^2 p_n - \log p_n\) \textbf{is} monotone in the
index; \(f_n\) is not. Substituting one for the other is the mistake,
and it is what makes FFM "the price of exactness".

Tier \textbf{L0 + {[}C{]}}.

\subsubsection{6.3 Unconditional analytic extension of the verified
range --- dead
permanently}\label{63-unconditional-analytic-extension-of-the-verified-range--dead-permanently}

\begin{quote}
\textbf{Theorem F (established).} No unconditional analytic result in
the literature extends the verified range by a single prime, and none
can: the best proven gap bound overshoots the barrier by a factor that
\textbf{diverges} --- \(1.5\times10^{7}\) at \(10^{20}\),
\(6\times10^{47}\) at \(10^{100}\).
\end{quote}

The range is \textbf{data-bound, not method-bound}. Tier \textbf{L0 +
{[}C{]}}.

\begin{center}\rule{0.5\linewidth}{0.5pt}\end{center}

\subsection{7. Computational evidence}\label{7-computational-evidence}

Every number here is reproducible from the run directory; no RNG, no
network.

\subsubsection{7.1 What was computed, and on what
tier}\label{71-what-was-computed-and-on-what-tier}

\begin{longtable}[]{@{}lll@{}}
\toprule\noalign{}
range & certified by & tier \\
\midrule\noalign{}
\endhead
\bottomrule\noalign{}
\endlastfoot
\([2,\,1223]\) & exact integer comparison \(p_{n+1}^n < p_n^{n+1}\) &
\textbf{L0} --- no arithmetic model at all \\
\([2,\,10^9)\) & exhaustive sieve + runtime-certified error bound; 50
847 533 pairs, 9 s & \textbf{L0} --- owned outright \\
\([89,\,4\times10^{18})\) & lever (gap table + Axler) & \textbf{L2}, 71
inequalities \\
\([89,\,2^{64})\) & same lever, longer table & \textbf{L2}, 76
inequalities \\
\([89,\,10^{20})\) & same lever, full b-file & \textbf{L2/L3} (community
census), 80 inequalities \\
\([89,\,1.857\times10^{19})\) & second, independent lever needing no
Axler bound & \textbf{L2}, 79 inequalities \\
beyond \(10^{20}\) & nothing & open \\
\end{longtable}

\textbf{No counterexample was found at any scale examined.} No leg
expected to find one, and none of these ranges is evidence about the
conjecture.

\subsubsection{7.2 Independent
reproduction}\label{72-independent-reproduction}

An adversarial audit leg re-derived the load-bearing numbers in code
written independently of the producing legs. Selected results:

\begin{itemize}
\tightlist
\item
  Re-sieved to \(2\times10^{7}\) in independent NumPy:
  \(\min H_n = 1.196152422706632\) at \(n=2\);
  \(\max D_n = 0.5486599111467569\) at \(n=214\) (\(p=1307\)) ---
  \textbf{exact matches}.
\item
  Recomputed the record-locus bracket:
  \(1193.345367 \le 1193.417778 \le 1193.459723\) --- \textbf{reproduces
  to every quoted digit}.
\item
  Re-ran the 85-row lever: record \#4 fails
  (\(\underline{B}/G = 0.823\)), \#5 onward passes, tightest ratio
  \(1.054246\) at \#64 --- \textbf{matches}.
\item
  Cross-checked two independently sourced maximal-gap tables
  (Wikipedia-sourced against OEIS b-files): \textbf{85/85 rows
  identical} in \(G\), \(P\) and \(\pi(P)\). \emph{No producing leg had
  performed this cross-check.}
\item
  Tested the unconditional criterion empirically at every index to
  \(2\times10^{7}\): \textbf{0 violations}.
\item
  Re-ran all seven verifier scripts: \textbf{all exit 0}, four producing
  byte-identical output.
\end{itemize}

\subsubsection{7.3 A genuinely open computational
obligation}\label{73-a-genuinely-open-computational-obligation}

Monotonicity of the lower barrier bound \(\underline{B}\) is used and
\textbf{not proved}. It was tested on a grid \textasciitilde500× denser
than any producing leg used --- 300 001 points over \([10, 10^{22}]\)
--- with \textbf{0 non-monotone steps}. The gap remains a gap. Recorded
as an open obligation, not as a result.

\subsubsection{7.4 A correction to this work\textquotesingle s own
opening
framing}\label{74-a-correction-to-this-works-own-opening-framing}

The attack opened with "the refutation branch is the live branch; the
shortfall is \(5.5\%\)". That framing was \textbf{partially demoted}
during the run, and this paper carries the demotion rather than the
headline:

\begin{itemize}
\tightlist
\item
  the shortfall figure is arithmetically right (corrected: \(1.05426\),
  §5.1) but is \textbf{not L0} as originally derived --- that derivation
  traversed an unclosed literature node; it is L0 \emph{now}, via the
  published exact barrier;
\item
  the sufficient criterion the headline leans on is \textbf{silent at
  \(n = 2\) and \(n = 4\)} --- the \emph{tightest known cases} of the
  conjecture. F is true there by the exact predicate (\(25 < 27\);
  \(14641 < 16807\)), but the criterion does not know it;
\item
  of two \textbf{pre-registered} demotion criteria, one fired and one
  did not. The naive probabilistic model predicts \(6.72\)
  counterexamples below the verified frontier where \(0\) are observed
  --- over-predicting by \(6.7\times\); the fitted crossover height is
  \(10^{30}\)--\(10^{51}\), not \(10^{21}\).
\end{itemize}

\textbf{The honest reading}: the refutation branch keeps its heuristic
motivation --- Granville\textquotesingle s
\(\limsup_{x} \max_{p_n \le x}(p_{n+1}-p_n)/\log^2 x \gtrsim 2e^{-\gamma} \approx 1.12292\)
is not excluded, and this \(\limsup\) is genuinely unknown --- and
\textbf{loses its quantitative one}. Search is not affordable at any of
the fitted heights.

The crossover-height figure itself carries a live caveat: it is a
\textbf{30-decade extrapolation from a model fit}, and the model tag did
not survive into the decision that used it. It is reported here as a
model output, never as a property of the primes.

\begin{center}\rule{0.5\linewidth}{0.5pt}\end{center}

\subsection{8. Evidence gate: BLOCKED}\label{8-evidence-gate-blocked}

The gate is fail-closed and it closed. This section reports the verdict
as the gate returned it.

\begin{longtable}[]{@{}lll@{}}
\toprule\noalign{}
leg & requirement & verdict \\
\midrule\noalign{}
\endhead
\bottomrule\noalign{}
\endlastfoot
\textbf{LOOP} & verdict file present, well-formed, names the live round
& \textbf{PASS} --- live round = 1 \\
\textbf{KERNEL} & \texttt{lake\ build} exit 0 \textbf{and} grep-clean of
\texttt{sorry}/\texttt{axiom} & \textbf{FAIL} --- exit 0 ✅ but one live
\texttt{sorry}; 7 declarations carry \texttt{sorryAx} \\
\textbf{SKEPTIC} & fault report exists \textbf{and} zero residual
BLOCKERs & \textbf{FAIL} --- 2 live BLOCKERs \\
\textbf{CORPUS} & red-team corpus present \textbf{and} coverage report
non-empty & \textbf{PASS} --- 19 entries / 15 categories \\
\end{longtable}

\textbf{FAIL, not DEGRADED.} DEGRADED is available only when no formal
backend was attempted. Lean \emph{was} run, \emph{did} build, and
\emph{did not} discharge the \texttt{sorry}. Calling that DEGRADED would
convert "we tried and the conjecture is still open" into "we could not
try" --- a different and false statement.

\subsubsection{8.1 The two live
BLOCKERs}\label{81-the-two-live-blockers}

The audit found \textbf{no mathematical error in any theorem}. Both
BLOCKERs are claims the corpus makes \emph{about its own
trustworthiness}, and both are false. That is exactly the class of fault
a gate exists to stop.

\begin{longtable}[]{@{}ll@{}}
\toprule\noalign{}
id & fault \\
\midrule\noalign{}
\endhead
\bottomrule\noalign{}
\endlastfoot
\textbf{BLOCKER-1} & A provenance record states that 44 rows of the
load-bearing maximal-gap table were independently recomputed.
\textbf{Thirty were.} A 47\% overstatement, sitting in the provenance
record of the corpus\textquotesingle s single weakest premise (P2 of
Theorem A). \\
\textbf{BLOCKER-2} & The claim \emph{"the verified range itself excludes
\(C \ge 1.1736\)"} is an \textbf{invalid inference}: a finite range
cannot exclude a value of a \(\limsup\). It is a property of a model
fit, read as a statement about the tail. \\
\end{longtable}

Neither needs new mathematics. One is a wrong row count; one is a
sentence that must be deleted or restated. \textbf{This paper does not
rely on either claim}: §7.2 reports only counts the audit itself
verified, and §7.4 states the \(\limsup\) as unknown.

\subsubsection{8.2 Seven MAJOR faults, carried
forward}\label{82-seven-major-faults-carried-forward}

Dominated by citation and cross-leg consistency: a locator absent from
the version of record (\textbf{MAJOR-6}, disclosed in §5.2); two
mutually incompatible verification tiers for the same source
(\textbf{MAJOR-7}); one leg crediting a sibling with a premise the
sibling had retracted (\textbf{MAJOR-4}); a vacuous conclusion sold as a
headline verdict (\textbf{MAJOR-3}, corrected in §5.3); a docstring
carrying a formally retracted claim (\textbf{MAJOR-5}); an undated
superlative survey claim inside a V0 row (\textbf{MAJOR-8}, see §11.2);
and the 30-decade extrapolation feeding a budget decision
(\textbf{MAJOR-9}, caveated in §7.4).

Of these, \textbf{MAJOR-3 and MAJOR-9 are corrected in this paper};
\textbf{MAJOR-6 and MAJOR-8 are disclosed at the point of use};
\textbf{MAJOR-4, MAJOR-5 and MAJOR-7 concern upstream artifacts this
paper does not rely on.}

\begin{center}\rule{0.5\linewidth}{0.5pt}\end{center}

\subsection{9. Limitations --- stated in full, not
summarised}\label{9-limitations--stated-in-full-not-summarised}

\begin{enumerate}
\def\labelenumi{\arabic{enumi}.}
\tightlist
\item
  \textbf{The conjecture is open.} It is not proved and not refuted
  here. Nothing in this paper is evidence about its truth.
\item
  \textbf{No finite range is evidence.} Neither the kernel-checked
  \(n \le 25\), nor the exhaustive \(p < 10^9\), nor Theorem
  A\textquotesingle s \(10^{20}\). The whole difficulty is in the tail.
\item
  \textbf{Theorem A inherits an unrefereed computational premise (P2,
  L3/V1)} and a \textbf{faulted citation locator (P1, MAJOR-6)}. It
  should be read as \emph{medium} confidence, with the weak link on the
  data side.
\item
  \textbf{The citation audit has not run.} No citation here has been
  independently re-fetched and re-read by an auditor. §11 traces each to
  a ledger row; that is provenance, not clearance.
\item
  \textbf{Six citations used in §6.1 sit on ledger rows that were
  proposed but never merged} into the authoritative ledger. §11.2 lists
  them explicitly.
\item
  \textbf{Three named gaps remain unsourced or unusable}: the four-term
  expansion \(B_n = L^2-L-1-3/L-13/L^2+O(L^{-3})\) is \textbf{unsourced}
  --- the published theorem gives only the \(+o(1)\) form, and this
  paper therefore does \textbf{not} use the lower-order terms; the
  Baker--Harman--Pintz threshold \(x_0\) is \textbf{ineffective as
  published}, so \(0.525\) must never be quoted as an explicit
  unconditional bound without that caveat; and no live URL exists for
  the canonical Nicely gap tables (all return HTTP 404 as of 2026-07-24)
  --- cite the community project instead.
\item
  \textbf{The monotonicity of \(\underline{B}\) (§7.3) is used and
  unproved.}
\item
  \textbf{Only one attack round ran} (a configured cap, not a failure
  --- the executor was verified capable). There is therefore no
  shrink-versus-churn signal: the list of \texttt{sorry}-carrying
  declarations entered at 7 and left at 7, \textbf{identical, not
  churned}. The six corollaries were never independent targets --- they
  inherit from the head.
\item
  \textbf{\texttt{lake\ build} was not re-executed by this leg.} The
  build record (exit 0, 8253 jobs) is read as the kernel
  leg\textquotesingle s own transcript.
\item
  \textbf{The delivery posture is staged.} §0.2.
\end{enumerate}

\subsubsection{9.1 What must close before
release}\label{91-what-must-close-before-release}

In the order the evidence imposes:

\begin{enumerate}
\def\labelenumi{\arabic{enumi}.}
\tightlist
\item
  Close \textbf{BLOCKER-1} (correct the row count) and
  \textbf{BLOCKER-2} (delete or restate the invalid inference), and
  unwind any decision resting on the latter.
\item
  Repair the \textbf{MAJOR-6} locator: establish where the \(3.83\)
  member sits in the \emph{published} Axler numbering, or re-attribute
  P1 to arXiv v3 explicitly in every place it appears.
\item
  Merge or drop the six \textbf{proposed-but-unmerged ledger rows} of
  §11.2. Theorems C--E cannot ship on unmerged rows.
\item
  Run \textbf{\texttt{citation-gate}}, then
  \textbf{\texttt{editorial-verdict}}.
\end{enumerate}

Steps 1--4 are corpus-honesty and citation work, not mathematics.
\textbf{None of them will close the kernel leg}: that closes only if the
conjecture is proved, which is the open problem itself.

\subsubsection{9.2 The honest forecast}\label{92-the-honest-forecast}

More rounds of the same shape would very likely keep producing
scaffolding without moving the head. The obstruction is not a missing
lemma that a re-read would surface --- it is the \textbf{categorical gap
between polylogarithmic and power-of-\(p\) gap bounds} (§1.4). If
further rounds are run, their value is in closing the corpus faults and
hardening the criteria, \textbf{not} in expecting the \texttt{sorry} to
fall.

\begin{center}\rule{0.5\linewidth}{0.5pt}\end{center}

\subsection{10. What this paper does not
claim}\label{10-what-this-paper-does-not-claim}

\begin{itemize}
\tightlist
\item
  It does not claim the conjecture is true, false, likely, or unlikely.
  The community\textquotesingle s heuristic expectation (from
  Granville\textquotesingle s revised Cramér model) is that it is
  \textbf{false}; that is recorded as a heuristic and is used as a
  premise nowhere.
\item
  It does not claim citation clearance (§9, limitation 4; §11).
\item
  It does not claim that RH fails to imply Firoozbakht (§6.1).
\item
  It does not claim any verified range is evidence.
\item
  It does not claim FFM on any range as a substantive result (§5.3).
\item
  It does not use the unsourced lower-order terms of the \(B_n\)
  expansion.
\item
  It does not claim "proved" for anything outside §4.
\end{itemize}

\begin{center}\rule{0.5\linewidth}{0.5pt}\end{center}

\subsection{11. Citation trace}\label{11-citation-trace}

Every citation in this paper resolves to a row in the
run\textquotesingle s authoritative source ledger,
\texttt{source-ledger/source-ledger.md}. Section numbers below are that
document\textquotesingle s.

\subsubsection{11.1 Merged ledger rows used in this
paper}\label{111-merged-ledger-rows-used-in-this-paper}

\begin{longtable}[]{@{}lllll@{}}
\toprule\noalign{}
citekey & ledger § & V & locator relied on & used in \\
\midrule\noalign{}
\endhead
\bottomrule\noalign{}
\endlastfoot
\texttt{firoozbakht1982} & §3.1 & V1 & attribution + date only & §1.1,
§0.3 \\
\texttt{rivera2002conj30} & §3.1 & V1 & page body & §1.1 \\
\texttt{ribenboim2004little} & §3.1 & V2 & p. 185 & §1.1 \\
\texttt{visser2019verifying} & §3.1, §3.2, §3.8 & \textbf{V0} & Conj. 3
eq. (2.4); Abstract; §4 & §1.3, §4.1, §7.1 \\
\texttt{kourbatov2015verification} & §3.2 & V1 & Abstract & §4.1,
§7.1 \\
\texttt{kourbatov2015upper} & §3.3, §3.4 & \textbf{V0 + C} & Table 1
last row; Thm 5; §3 & §1.3, §5.1 \\
\texttt{kourbatov2015corrigendum} & §3.4 & \textbf{V0} & full text &
§11.3 \\
\texttt{primegaplist\_faq} & §3.2, §3.3 & V1 + C & FAQ body & §5.1,
§5.2, §7.1 \\
\texttt{axler2014newbounds} & §3.5 & \textbf{V0} & Cor. 3.5 (arXiv v3)
--- \textbf{see MAJOR-6} & §5.2 \\
\texttt{axler2016corrigendum} & §3.5 & V2 & via Kourbatov corrigendum &
§5.2, §11.3 \\
\texttt{dusart2010estimates} & §3.5 & \textbf{V0 + C} & Props. 6.6, 6.7
& §5.1 \\
\texttt{granville1995harald} & §3.6 & \textbf{V0} & display following
eq. (20); eq. (14) & §7.4 \\
\texttt{cramer1936order} & §3.6 & V2 & block quotation in Granville &
§11.3 \\
\texttt{bhp2001difference} & §3.6 & V2 & main theorem & §1.4; §9,
limitation 6 \\
\texttt{fgkmt2018long} & §3.7 & V1 & Abstract & §11.3 \\
\end{longtable}

\subsubsection{\texorpdfstring{11.2 Proposed but \textbf{unmerged} rows
--- a live
obligation}{11.2 Proposed but unmerged rows --- a live obligation}}\label{112-proposed-but-unmerged-rows--a-live-obligation}

The RH-route results of §6.1 rest on citations that a downstream leg
\textbf{proposed} for the ledger and that were \textbf{never merged}
into it. They are recorded here with the verification level the
proposing leg claimed, and they are \textbf{not cleared}:

\begin{longtable}[]{@{}llll@{}}
\toprule\noalign{}
citekey & claimed V & statement relied on & status \\
\midrule\noalign{}
\endhead
\bottomrule\noalign{}
\endlastfoot
\texttt{cms2019fourier} --- Carneiro, Milinovich, Soundararajan,
\emph{Fourier optimization and prime gaps}, Comment. Math. Helv.
\textbf{94} (2019), 533--568; arXiv:1708.04122 & \textbf{V0} (claimed;
quotations independently re-fetched and confirmed exact by the audit
leg) & on RH, \(g_n \le \tfrac{22}{25}\sqrt{p_n}\log p_n\) for
\(p_n > 3\); the method\textquotesingle s limiting constant is
\(\tfrac12\) & \textbf{UNMERGED} \\
\texttt{littlewood1914} & \textbf{V3} &
\(\psi(x)-x = \Omega_\pm(\sqrt x \log\log\log x)\) & \textbf{UNMERGED +
V3} --- used only for the explanatory remark in §6.1, which is flagged
as such \\
\texttt{dudek2015riemann} & V1 & earlier RH-conditional constant &
\textbf{UNMERGED}, lineage only \\
\texttt{leenosal2025sharper} & \textbf{V0} & sharper RH error bounds &
\textbf{UNMERGED}, not load-bearing here \\
\texttt{schoenfeld1976sharper} & V3 → V2 (tier disputed,
\textbf{MAJOR-7}) & classical RH bound & \textbf{UNMERGED}, \textbf{not
used in this paper} \\
\texttt{montgomery1973pair} & V3 & pair-correlation consequence &
\textbf{UNMERGED}, \textbf{not used in this paper} \\
\end{longtable}

\textbf{Consequence, stated plainly.} Theorems C, D and E of §6.1 are
\textbf{established} mathematics resting on a citation that is not yet
in the authoritative ledger. They must not be released until
\texttt{cms2019fourier} is merged at V0 and \texttt{littlewood1914} is
either promoted or the explanatory remark is restated as conditional on
it. This is item 3 of §9.1.

A further caveat on the same family (\textbf{MAJOR-8}): the phrase
\emph{"the strongest published RH-conditional gap bound"} is a
\textbf{superlative survey claim}, and no leg documents a literature
sweep after 2019 --- one was attempted and blocked by rate limiting.
§6.1 therefore says "the strongest \emph{published} bound" as the
producing leg\textquotesingle s assessment, not as a verified fact about
the literature as of 2026.

\subsubsection{11.3 Citation hygiene notes that must survive into any
restatement}\label{113-citation-hygiene-notes-that-must-survive-into-any-restatement}

Five errors are easy to make in this literature, and the ledger pins
each:

\begin{enumerate}
\def\labelenumi{\arabic{enumi}.}
\tightlist
\item
  \textbf{The \(b = 1.17\) sufficient condition has a corrected validity
  range.} Kourbatov\textquotesingle s Theorem 3 originally rested on
  Axler\textquotesingle s bound at "\(x \ge 5.43\)";
  Axler\textquotesingle s own corrigendum moved that to
  \(x \ge 2\,634\,800\,823\), and Kourbatov published a corrigendum
  rewriting the proof. Any citation of "\(b = 1.17 \Rightarrow\)
  Firoozbakht" that does not carry \(p_k \ge 2\,634\,800\,823\)
  \textbf{plus} the unconditional check on \([29,\, 2\,634\,800\,823]\)
  is citing a withdrawn proof. Cite arXiv:1506.03042\textbf{v4}.
\item
  \textbf{The verification frontier is \(2^{64}\), not
  \(4\times10^{18}\).} \(4\times10^{18}\) is Kourbatov\textquotesingle s
  2015 exhaustive-sieve range; Visser 2019 covers \(p < 2^{64} \approx
  1.844\times10^{19}\), for Firoozbakht \emph{and} the two strictly
  stronger Nicholson and Farhadian conjectures. The community census
  reaches \(10^{20}\) (§5.2, P2).
\item
  \textbf{Explicit bounds here are Axler\textquotesingle s, not
  Rosser--Schoenfeld\textquotesingle s or Dusart\textquotesingle s.}
  Kourbatov\textquotesingle s chain uses Axler. Writing
  "Rosser--Schoenfeld / Dusart" while reproducing that derivation is
  mis-citing. (Dusart 2010 \emph{is} used, at V0, but for the \(p_k\)
  bracket of §5.1 --- a different role.)
\item
  \textbf{The long-gaps exponent is \((\log\log\log X)^{1}\), not
  squared.} The 2018 record improves the 2016 one by exactly one
  iterated log in that denominator. This is the easiest citation error
  in the family. Either way the bound is \(\ll \log^2 X\) --- too small
  by a factor \(\sim \log X/\log\log X\) to touch Firoozbakht.
\item
  \textbf{Cramér stated the \(\limsup = 1\) relation about his random
  model, and only \emph{suggested} the transfer to the primes.} Any
  sentence beginning "Cramér conjectured" must carry that hedge. And
  "Cramér 1920/21" is a mis-citation: the RH gap bound is Cramér 1920,
  with a distinct companion paper the same year.
\end{enumerate}

Also recorded as \textbf{absences}, which is part of the provenance:
there is \textbf{no primary document by Firoozbakht} (the standard
reference is literally "\emph{(1982), unpublished}"), so "Firoozbakht
proposed, in 1982" is supportable and "Firoozbakht proved/published" is
not; the canonical Nicely gap tables are \textbf{dead links}; and
\textbf{no proof, disproof, or accepted counterexample of the conjecture
was found anywhere in the literature.}

The bibliography accompanying this paper is \texttt{references.bib} (31
merged citekeys, plus a clearly separated addendum block for
§11.2\textquotesingle s unmerged rows).

\begin{center}\rule{0.5\linewidth}{0.5pt}\end{center}

\subsection{12. Reproduction}\label{12-reproduction}

All artifacts of the run live under the run directory. The load-bearing
checks:

\begin{verbatim}
# formal layer
cd <run>/lean-probe/lean && lake build          # exit 0; one `sorry` warning, at Statement.lean
                                                # Audit.lean prints axioms for 31 declarations

# recomputation layer (each deterministic, no RNG, no network)
python3 <run>/source-ledger/verify_ledger.py    # ledger rows
python3 <run>/proof-attempt__0/verify_pa0.py    # FFM numerics
python3 <run>/proof-attempt__1/verify_rh.py     # RH usable sets
python3 <run>/proof-attempt__2/verify_pa2.py    # the 85-row lever
python3 <run>/notebooks__2/verify-findings.py   # the verified-range lever
python3 <run>/concept-cards/verify_cards.py     # concept-card invariants
python3 <run>/skeptic/verify_faults.py          # the audit's own re-derivations
python3 <run>/synthesize/verify_synthesis.py    # 56/56 checks, exit 0
\end{verbatim}

The synthesis verifier is deliberately a \textbf{tripwire}: it asserts
that the two BLOCKERs are still live. If a future session closes them,
those checks \textbf{fail} --- which is the intended behaviour, not a
regression.

\begin{center}\rule{0.5\linewidth}{0.5pt}\end{center}

\subsection{13. The one-paragraph
version}\label{13-the-one-paragraph-version}

The primes walk along a road, and every step is a gap. Firoozbakht drew
a fence at height roughly \((\log p)^2\) and said the primes never jump
it --- not once, not ever, all the way to infinity. This work did not
decide that. What it did was survey the ground carefully. It showed the
fence is \emph{exactly} the right description of the conjecture, not an
approximation. It proved --- inside a kernel, in the literal form ---
that the primes stay under it for the first 25 steps, and established by
argument that they do so for every prime below \(10^{20}\), resting on
two named external premises. It showed that the Riemann hypothesis, the
biggest hammer anyone would reach for, settles the question for the
prime 5 and for no other. It showed that no unconditional method in the
literature can extend the verified range by a single prime, and that the
shortfall \textbf{diverges} rather than closing. And it built the two
exact criteria that a future proof and a future refutation would each
need. Then its own auditor found two places where the corpus was
overstating how carefully it had checked itself, and the gate closed red
on that. \textbf{The mathematics in this corpus is sound. Its
self-description was not, in two places. The conjecture is still open.}

\begin{center}\rule{0.5\linewidth}{0.5pt}\end{center}

\emph{Emitted by the \texttt{write-paper} leg
(\texttt{edit-20260724-9e06}) of the \texttt{math-attack} polymer, run
node. External attribution: \textbf{Noogram}.
Firoozbakht\textquotesingle s conjecture is \textbf{open} --- not
proved, not refuted. Evidence-gate status: \textbf{BLOCKED}. Citation
clearance: \textbf{not claimed, not audited}. Delivery posture:
\textbf{staged}.}

\printbibliography

\end{document}
